\documentclass[11pt,a4paper]{article}

\usepackage{../shared/cathedral-preamble}
\usepackage{setspace}
\onehalfspacing

\title{%
  Who Owns a Machine-Verified Proof?\\
  Legal and Intellectual Property Questions\\
  from the Cathedral Project%
}
\author{
  Jason Robert Gochanour
}
\date{April 20, 2026}

\begin{document}

\maketitle

\begin{abstract}
The Cathedral---a machine-verified reduction of the Riemann
Hypothesis---was built by a human with AI assistance. This paper
examines the legal and intellectual property questions raised by
this mode of production: (1)~copyright and ownership of
AI-assisted mathematical proofs, (2)~the legal status of AI as
co-author or co-inventor, (3)~machine-verified proofs as legal
evidence, (4)~open-source licensing of formal mathematics, and
(5)~liability for AI-generated mathematical errors. We do not
provide legal advice. We identify the questions that existing law
does not clearly answer and suggest that these questions will
become urgent as AI-assisted formal verification becomes
commonplace. This paper is written for attorneys, IP professionals,
and policymakers who may encounter these issues in practice.
\end{abstract}

\tableofcontents

% ================================================================
\section{The Facts of the Case}
% ================================================================

The Cathedral was produced under the following circumstances:

\begin{itemize}
  \item \textbf{Human contributor}: One individual (the author),
    a computer scientist, providing architectural design,
    mathematical strategy, quality control, and editorial judgment.

  \item \textbf{AI contributors}: Two commercial AI systems
    (Google DeepMind's Gemini and Anthropic's Claude) providing
    code generation, proof strategy suggestions, and mathematical
    analysis. Both were used through standard commercial APIs and
    subscriptions.

  \item \textbf{Verification}: The Lean~4 theorem prover, an
    open-source software system, verified every logical step.
    The Mathlib library (open-source, Apache~2.0 license) provided
    foundational mathematics.

  \item \textbf{Output}: 78 Lean source files, 96 archived files,
    and 22 companion papers. Approximately 36,738 lines of
    verified proof code and 60,000 words of documentation.

  \item \textbf{Timeline}: 23 days of development.
\end{itemize}

% ================================================================
\section{Question 1: Who Owns the Proof?}
\label{sec:ownership}
% ================================================================

\subsection{Copyright in Mathematical Proofs}

Under U.S. copyright law (17 U.S.C. \S~102), copyright protects
``original works of authorship fixed in any tangible medium of
expression.'' Mathematical proofs have an ambiguous status:

\begin{itemize}
  \item The \emph{mathematical content} (theorems, identities,
    logical relationships) is not copyrightable. Mathematical
    facts are discoveries, not inventions.
  \item The \emph{expression} of a proof (the specific sequence
    of words, the pedagogical choices, the narrative structure)
    is copyrightable.
  \item \emph{Computer code} implementing a proof is copyrightable
    as a literary work.
\end{itemize}

The Cathedral's Lean code is original expression implementing
mathematical content. The code itself is copyrightable; the
mathematical results it proves are not.

\subsection{AI-Generated Content}

The U.S. Copyright Office has consistently held that copyright
requires a human author. In \emph{Th\'aler v. Perlmutter} (2023),
the D.C. District Court upheld the Copyright Office's refusal to
register a work created autonomously by AI.

However, the Cathedral is not autonomously AI-generated. The human
provided:
\begin{itemize}
  \item The overall architecture (which modules exist, how they
    connect).
  \item The mathematical strategy (Parseval Bridge over Vasyunin,
    Rank-1 Mellin over Hahn--Banach).
  \item Quality control (accepting, rejecting, or modifying
    AI-generated code).
  \item Editorial judgment (what to archive, what to keep, how
    to document).
\end{itemize}

This is analogous to a photographer using a camera: the camera
captures the image, but the photographer makes the creative
decisions (composition, timing, subject). Courts have consistently
held that photographs are copyrightable despite the mechanical
role of the camera.

\textbf{Likely outcome}: The Cathedral's code is copyrightable by
the human author, as the human provided sufficient creative
direction. The AI's contribution is analogous to a camera's or a
compiler's: mechanical execution of human-directed creative choices.

\subsection{Terms of Service Considerations}

Both Google and Anthropic's terms of service (as of 2026) assign
ownership of AI-generated output to the user. However:

\begin{itemize}
  \item These terms could change.
  \item They may not be enforceable in all jurisdictions.
  \item They do not address the underlying copyright question
    (whether the output is copyrightable at all).
\end{itemize}

% ================================================================
\section{Question 2: Can AI Be a Co-Author?}
\label{sec:authorship}
% ================================================================

\subsection{Academic Authorship}

The Cathedral's companion papers credit AI systems as
``collaborative contributors.'' This raises questions under
academic authorship standards:

\begin{itemize}
  \item \textbf{ICMJE criteria}: Authorship requires (a)~substantial
    contribution, (b)~drafting or critical revision, (c)~final
    approval, and (d)~accountability. AI systems cannot satisfy
    (c) or (d).
  \item \textbf{ACM policy}: The ACM prohibits listing AI systems
    as authors but requires disclosure of AI assistance.
  \item \textbf{Nature/Science}: Major journals require that all
    authors be human and accountable.
\end{itemize}

The Cathedral's documentation credits AI as a contributor in the
acknowledgments, not in the author line. This is consistent with
current academic norms but may understate the AI's contribution
(85\% of code by volume, strategy proposals adopted by the human).

\subsection{Inventorship}

Under U.S. patent law (35 U.S.C. \S~100), an inventor must be a
``natural person.'' In \emph{Th\'aler v. Vidal} (2022), the
Federal Circuit held that AI cannot be named as an inventor on a
patent.

The Cathedral's methods (Parseval Bridge, Rank-1 Mellin Miracle)
are mathematical techniques, not patentable inventions. However,
the engineering applications proposed in companion papers (M\"obius
Harmonic Analyzer, arithmetic compressed sensing, prime-spaced
arrays) could potentially be patentable. If AI proposed these
applications, the inventorship question becomes relevant.

\textbf{Current law}: Only the human can be named as inventor,
even if AI substantially contributed to the inventive concept.
This may change as legislatures address AI inventorship.

% ================================================================
\section{Question 3: Machine-Verified Proofs as Legal Evidence}
\label{sec:evidence}
% ================================================================

\subsection{The Evidentiary Question}

Machine-verified proofs could serve as evidence in legal
proceedings:

\begin{itemize}
  \item \textbf{Patent disputes}: ``Our algorithm is provably
    correct'' (verified in Lean~4).
  \item \textbf{Safety certification}: ``This bridge design
    meets load requirements'' (verified by formal analysis).
  \item \textbf{Regulatory compliance}: ``Our financial model
    satisfies Basel III constraints'' (verified).
  \item \textbf{Product liability}: ``The software was formally
    verified to be free of the defect that caused the accident.''
\end{itemize}

\subsection{Admissibility}

Under the Federal Rules of Evidence (FRE 702), expert testimony
must be based on ``sufficient facts or data'' and ``reliable
principles and methods.'' A machine-verified proof could qualify
as:

\begin{itemize}
  \item \textbf{Expert testimony}: The Lean compiler's output as
    an expert opinion on the correctness of a mathematical
    argument.
  \item \textbf{Documentary evidence}: The Lean source code as a
    document establishing a mathematical fact.
  \item \textbf{Computer-generated evidence}: Similar to DNA
    analysis or financial modeling output, subject to
    authentication (FRE 901) and the Daubert standard.
\end{itemize}

\textbf{Key advantage}: Unlike human expert testimony, machine-
verified proofs are \emph{deterministic} and \emph{reproducible}.
Any party can re-run the Lean compiler and obtain the same result.
This satisfies the Daubert factors (testable, peer-reviewed,
known error rate [zero], generally accepted in the relevant
community).

\textbf{Key challenge}: The proof's correctness depends on its
\emph{axioms}. An opposing party could challenge the axioms rather
than the proof. The Cathedral model---where every axiom is named,
classified, and documented---facilitates this: the debate shifts
from ``is the proof correct?'' (machine-checked) to ``are the
assumptions reasonable?'' (human judgment).

% ================================================================
\section{Question 4: Open-Source Licensing}
\label{sec:license}
% ================================================================

\subsection{The Licensing Landscape}

The Cathedral uses and depends on open-source components:

\begin{center}
\begin{tabular}{@{}lll@{}}
\toprule
\textbf{Component} & \textbf{License} & \textbf{Implication} \\
\midrule
Lean 4 & Apache 2.0 & Permissive, commercial use OK \\
Mathlib & Apache 2.0 & Permissive, commercial use OK \\
Cathedral code & Not yet licensed & See below \\
Companion papers & Not yet licensed & See below \\
\bottomrule
\end{tabular}
\end{center}

\subsection{Licensing Considerations}

The Cathedral's author must decide on a license. Options include:

\begin{enumerate}
  \item \textbf{Apache 2.0 / MIT}: Consistent with Lean/Mathlib
    ecosystem. Allows commercial use and modification. Maximizes
    adoption. Does not prevent military or surveillance use.

  \item \textbf{GPL}: Requires derivative works to be open-source.
    Limits commercial exploitation. Inconsistent with Mathlib's
    Apache 2.0 license (potential compatibility issue).

  \item \textbf{Creative Commons (for papers)}: CC-BY for maximum
    dissemination. CC-BY-NC to prevent commercial exploitation
    of the documentation.

  \item \textbf{Custom license with ethical restrictions}: Could
    prohibit use for specific purposes (e.g., surveillance,
    weapons). Difficult to enforce. Not OSI-approved. May limit
    adoption.

  \item \textbf{No license (all rights reserved)}: Default under
    copyright law. Maximizes author control. Minimizes adoption
    and community contribution.
\end{enumerate}

\textbf{Recommendation}: Apache 2.0 for code (ecosystem
consistency), CC-BY-4.0 for papers (maximum dissemination with
attribution). This is consistent with the project's stated values
of transparency and open science, and with the companion papers'
argument that open publication favors defenders over attackers.

% ================================================================
\section{Question 5: Liability}
\label{sec:liability}
% ================================================================

\subsection{If the Proof Has an Error}

The Cathedral claims zero compilation errors and two explicit
axioms. But what if:

\begin{enumerate}
  \item \textbf{Lean has a bug}: The Lean compiler itself could
    have a soundness bug that accepts an invalid proof. This is
    theoretically possible but highly unlikely (Lean's kernel is
    small and well-audited).

  \item \textbf{An axiom is false}: One of the 47 axioms could be
    mathematically false. The proof would be ``correct'' conditional
    on a false assumption. The tiered axiom classification mitigates
    this: Tier~1 axioms are well-established results; Tier~4 axioms
    are structural claims that could potentially be wrong.

  \item \textbf{A formalization error}: The Lean code could
    correctly implement a theorem that does not accurately represent
    the intended mathematical statement. This is the ``specification
    bug'' problem.
\end{enumerate}

\subsection{Liability Framework}

If someone relies on a Cathedral-derived result (e.g., a verified
stability certificate for a bridge) and that result is wrong:

\begin{itemize}
  \item \textbf{Product liability}: Does the proof ``product''
    create liability for the author? Under current law, probably
    not---mathematical proofs are information, not products. But
    if the proof is packaged as a commercial service (``certified
    correct by formal verification''), product liability could
    attach.

  \item \textbf{Professional liability}: If the author holds
    themselves out as providing engineering verification services,
    professional liability standards apply.

  \item \textbf{AI liability}: If the AI contributed to an error,
    who is liable---the AI provider (Google, Anthropic), the user,
    or the compiler? Current law does not clearly answer this.
\end{itemize}

\textbf{Mitigation}: The Cathedral's axiom transparency is a
liability shield. Every assumption is documented. A user who relies
on the proof accepts the axioms; if an axiom is false, the proof's
warranty is limited to ``correct conditional on stated axioms.''
This is analogous to an engineering report that states its
assumptions explicitly.

% ================================================================
\section{Emerging Legal Questions}
% ================================================================

The Cathedral anticipates several legal questions that will become
urgent as AI-assisted formal verification becomes common:

\begin{enumerate}
  \item \textbf{Verification as due diligence}: Will formal
    verification become a legal standard of care? If a bridge
    \emph{could} have been formally verified but wasn't, does
    the engineer bear liability for using a less rigorous method?

  \item \textbf{Axiom disclosure}: If a verified claim depends on
    axioms, should there be a legal requirement to disclose them?
    (Analogous to financial disclosure requirements.)

  \item \textbf{AI training on proofs}: If an AI system is trained
    on the Cathedral's code, does the Cathedral's license apply to
    the AI's output? (The ``derivative work'' question for AI
    training data.)

  \item \textbf{Jurisdictional issues}: Mathematics is universal,
    but IP law is jurisdictional. A proof verified in the U.S.
    has no special legal status in the EU or China.

  \item \textbf{National security}: Can a government classify a
    mathematical proof? The dual-use paper identifies verified
    instability certificates as potentially sensitive information.
    Current frameworks (CEII, ITAR) do not clearly cover
    mathematical proofs.
\end{enumerate}

% ================================================================
\section{Conclusion}
% ================================================================

The Cathedral raises more legal questions than it answers. This is
expected: the law follows technology, not the reverse. The key
insight for legal professionals is that machine-verified proofs
create a \emph{new category} of intellectual output:

\begin{itemize}
  \item More rigorous than expert opinion (machine-checked).
  \item More transparent than traditional proof (axioms named).
  \item More reproducible than simulation (deterministic compilation).
  \item Created by a novel collaboration (human + AI + compiler).
\end{itemize}

Existing legal frameworks---copyright, patent, evidence,
licensing, liability---were designed for a world where intellectual
work is produced by humans using passive tools. The Cathedral is
produced by a human using active tools (AI) verified by an
objective tool (compiler). The legal frameworks will need to adapt.

The recommendation to legal professionals: study the Cathedral's
methodology now, before the first case arrives on your desk.
The questions are coming. They are inevitable. The only question
is whether the law is prepared.

\begin{quote}
\emph{The compiler does not practice law.\\
But it produces evidence that lawyers\\
will need to understand.}
\end{quote}

\begin{thebibliography}{99}

\bibitem{thaler2023}
S.~Thaler v. R.~Perlmutter, No. 1:22-cv-01564 (D.D.C. Aug.~18,
2023).

\bibitem{thaler2022}
S.~Thaler v. K.~Vidal, No. 21-2347 (Fed. Cir. Aug.~5, 2022).

\bibitem{copyright2023}
U.S. Copyright Office, \emph{Copyright Registration Guidance:
Works Containing Material Generated by Artificial Intelligence},
88 Fed. Reg. 16190 (Mar.~16, 2023).

\bibitem{daubert}
Daubert v. Merrell Dow Pharmaceuticals, 509 U.S. 579 (1993).

\bibitem{acm2023}
ACM, \emph{Policy on Authorship}, updated 2023.

\end{thebibliography}

\end{document}
